% Image credits for the photographs and AI illustrations of Book 5
% (University Biology, Year 3). Machine-readable license records live in
% images/book5/CREDITS.md; keep the two files in sync.
\chapter*{Créditos das imagens}

% \sloppy must be in force when the paragraph is BROKEN, hence the \par before
% \endgroup -- without it the group closes first and the tolerance is restored.
\begingroup\sloppy

Os desenhos deste livro são originais. As fotografias são usadas sob
as respectivas licenças livres, com agradecimento aos seus autores:

\begin{itemize}
  \item Capítulo~\ref{ch:b3:genomics}: fotografia de Frederick Sanger
        --- domínio público (Wikimedia Commons).
  \item Capítulo~\ref{ch:b3:genetic-engineering}: arroz dourado ao lado
        de arroz comum, International Rice Research Institute, 2011 ---
        CC~BY~2.0 (Wikimedia Commons).
  \item Capítulo~\ref{ch:b3:structural-biology}: fotografia de Max
        Perutz, Associated Press, 1962 --- domínio público (Wikimedia
        Commons).
  \item Capítulo~\ref{ch:b3:bacteriology}: retrato de Robert Koch,
        Wellcome Collection --- CC~BY~4.0 (Wikimedia Commons);
        fotografia de Alexander Fleming em seu laboratório, arquivo do
        U.S. Navy Bureau of Medicine and Surgery --- sem restrições de
        direitos conhecidas (Wikimedia Commons).
  \item Capítulo~\ref{ch:b3:innate-immunity}: fotografia de Élie
        Metchnikoff, Bain News Service, Library of Congress --- domínio
        público (Wikimedia Commons).
  \item Capítulo~\ref{ch:b3:adaptive-immunity}: retrato gravado de
        Edward Jenner, 1807 --- domínio público (Wikimedia Commons).
  \item Capítulo~\ref{ch:b3:nervous-systems}: desenho de uma célula de
        Purkinje por Santiago Ramón y Cajal, 1899 --- domínio público
        (Wikimedia Commons).
  \item Capítulo~\ref{ch:b3:endocrinology}: fotografia de Frederick
        Banting e Charles Best, c.~1924 --- domínio público (Wikimedia
        Commons).
  \item Capítulo~\ref{ch:b3:molecular-evolution}: fotografia de Motoo
        Kimura, 1986, fotógrafo desconhecido, via S.~Kumar e
        T.~Gojobori, \emph{Molecular Biology and Evolution} (2023) ---
        CC~BY~4.0 (Wikimedia Commons).
  \item Capítulo~\ref{ch:b3:behavioural-ecology}: fotografia de Niko
        Tinbergen por Rob Mieremet, Anefo, 1973 --- CC0, Nationaal
        Archief (Wikimedia Commons).
  \item Capítulo~\ref{ch:b3:conservation-biology}: fotografia de
        Martha, o último pombo-passageiro, Enno Meyer, 1912 --- domínio
        público (Wikimedia Commons); lobos no Parque Nacional de
        Yellowstone, National Park Service, 1996 --- domínio público
        (Wikimedia Commons); Sirocco, o kakapo, Department of
        Conservation, Nova Zelândia, 2012 --- CC~BY~2.0 (Wikimedia
        Commons).
\end{itemize}

Os links para as fontes completas e os endereços das licenças de cada
fotografia estão em
\texttt{images/\allowbreak book5/\allowbreak CREDITS.md}, no repositório do livro.

\bigskip
Ao lado das fotografias e dos desenhos originais, cerca de noventa
figuras espalhadas pelos capítulos do volume são ilustrações geradas
por inteligência artificial, produzidas com a geração de imagens da
OpenAI a partir de instruções escritas pelos editores do livro, e
revisadas quanto à exatidão biológica antes da inclusão. A instrução
completa de cada imagem gerada está registrada em
\texttt{images/\allowbreak book5/\allowbreak ai/\allowbreak PROMPTS.md}, no repositório.
\par\endgroup
\clearpage
